\mychapter{Materiais e Métodos}{cap:metodos}
\lhead{MATERIAIS E MÉTODOS}

A pesquisa é aplicada e experimental. Duas aplicações de referência, com telas de lista, formulário e navegação, serão projetadas no Figma com as bibliotecas oficiais de componentes de cada sistema operacional e implementadas em três versões: (A) escrita à mão em SwiftUI (macOS) e WinUI~3 (Windows); (B) gerada pelo framework a partir de uma fonte única; (C) com renderização própria, usando Electron como \textit{baseline}, um dos frameworks avaliados no desktop por Milichiewicz e Badurowicz \cite{milichiewicz2025comparative}. As três versões partem do mesmo projeto no Figma.

O framework será desenvolvido de forma iterativa. A fonte única será escrita em Kotlin, com sintaxe declarativa, e o gerador, preferencialmente em Rust, a lerá por análise sintática (parser externo ou plugin do compilador Kotlin, a definir) e produzirá os projetos nativos. O escopo se restringe a aplicações simples para Windows e macOS, e o Linux será tratado como trabalho futuro.

As medidas de desempenho são uso de CPU, memória e tempo de resposta, como em Biørn-Hansen et al.\ \cite{biorn2020empirical}, mais tamanho de armazenamento, como em Mahmoud et al.\ \cite{mahmoud2024trans}. Cada medição será repetida o suficiente para reportar mediana, intervalo de confiança e tamanho de efeito, e a H1 será avaliada por um teste de equivalência (como o TOST), com a margem fixada antes da coleta. A fidelidade visual será medida por similaridade estrutural (SSIM) entre capturas de tela das versões geradas e das escritas à mão, métrica também usada no GUIMigrator \cite{gao2026guimigrator}, e a cobertura de componentes nativos, por lista de conferência.

Os experimentos rodarão em uma máquina com macOS e em outra com Windows 11, ambas com hardware e versão do sistema mantidos fixos durante toda a coleta. Cada sistema será medido com as ferramentas de \textit{profiling} da própria plataforma, e as três versões rodarão nas mesmas condições (mesmo equipamento, mesma carga e estado inicial limpo). Todo o código, as telas de referência e os dados coletados ficarão no repositório do trabalho, para que a comparação possa ser replicada.
